\chapter{Localization}
\label{ch:t-localization}

\section{No map: the start spot and the tags}

The X30 patrols without a map. The robot starts on its dock; FAST-LIO's start is the origin of
the \code{map} frame, and every route is taught and repeated in that frame. Odometry drifts with
every metre; \textbf{tags} (printed ArUco markers on walls and posts, recorded once in the site)
correct the drift wherever the robot sees one.

\code{localization_node} estimates $\T{map}{odom}$, the correction between the map and the
odometry frame, with an extended Kalman filter on $SE(2)$:

\begin{center}
\includegraphics[width=\linewidth]{ekf}
\end{center}

\section{The filter}

\subsection{State and prediction}

The state is $\vect{x} = (x, y, \psi)$ of $\T{map}{odom}$ with its covariance $\mat{P}$. Between
corrections, odometry carries the robot and $\vect{x}$ stays; only the uncertainty grows, and
only with motion. Motion is counted in steps (\param{motion_step_m} \SI{5}{cm},
\param{motion_step_deg} \ang{2}) so that the sway of a standing robot adds nothing. For a step of
$\Delta s$ metres and $\Delta\theta$ radians, the process noise in the robot frame is
\begin{equation}
  \mat{Q}_{\text{base}} = \operatorname{diag}\!\left(\sigma_{xy}^2\,\Delta s,\ \sigma_{xy}^2\,\Delta s,\
    \sigma_{\psi,m}^2\,\Delta s + \sigma_{\psi,\text{rad}}^2\,|\Delta\theta|\right)
\end{equation}
mapped into the state through the inverse of the measurement Jacobian, $\mat{P} \leftarrow
\mat{P} + \mat{H}^{-1}\mat{Q}_{\text{base}}\mat{H}^{-\top}$. With the defaults
($\sigma_{xy} = 0.03\,\mathrm{m}/\sqrt{\mathrm{m}}$, $\sigma_{\psi,m} = 0.1^\circ/\sqrt{\mathrm{m}}$) a
\SI{50}{m} walk without a tag leaves $\sigma_{xy} = 0.03\sqrt{50} = \SI{0.21}{m}$ and
$\sigma_\psi = \ang{0.71}$.

\subsection{Correction}

A tag seen from the robot gives a measurement $\vect{z}$ of $\T{map}{base}$ with covariance
$\mat{R}$. The innovation against the prediction $h(\vect{x}) = \T{map}{odom}\,\T{odom}{base}$,
its covariance and the squared Mahalanobis distance are
\begin{equation}
  \vect{\nu} = \vect{z} \ominus h(\vect{x}),\qquad
  \mat{S} = \mat{H}\mat{P}\mat{H}^\top + \mat{R},\qquad
  d^2 = \vect{\nu}^\top \mat{S}^{-1} \vect{\nu}
\end{equation}
A measurement with $d^2$ above the $\chi^2_3$ gate is rejected as an outlier (a tag recorded
wrong, a tag moved, a false detection). Otherwise
\begin{equation}
  \mat{K} = \mat{P}\mat{H}^\top\mat{S}^{-1},\qquad
  \vect{x} \leftarrow \vect{x} + \mat{K}\vect{\nu},\qquad
  \mat{P} \leftarrow (\mat{I}-\mat{K}\mat{H})\,\mat{P}\,(\mat{I}-\mat{K}\mat{H})^\top +
    \mat{K}\mat{R}\mat{K}^\top
\end{equation}
The Joseph form keeps $\mat{P}$ symmetric and positive definite. After the update the
uncertainty is floored at \param{min_sigma_xy} (\SI{1}{cm}) and \param{min_sigma_yaw_deg}
(\ang{0.25}): the filter never becomes so sure of itself that it rejects the next tag.

\begin{center}
\begin{small}
\begin{tabular}{@{}lrrr@{}}
  \toprule
  \param{gate_confidence} & 0.95 & \textbf{0.99} & 0.999 \\
  $\chi^2_3$ gate & 7.81 & \textbf{11.34} & 16.27 \\
  \bottomrule
\end{tabular}
\end{small}
\end{center}

\subsection{How much a tag is trusted}

A tag's pose is less certain the farther and the more obliquely it is seen. For a tag at distance
$d$ seen at angle $\theta$ off its normal:
\begin{equation}
  \sigma_p = \sigma_{p,1}\,\frac{d^2}{\cos\theta},\qquad
  \sigma_\psi = \sigma_{\psi,1}\,\frac{d}{\cos\theta}
\end{equation}
with $\sigma_{p,1}$ \SI{2}{cm} and $\sigma_{\psi,1}$ \ang{1} at one metre
(\param{sigma_xy_at_1m}, \param{sigma_yaw_deg_at_1m}). A heading error also moves the robot about
the tag, so $\mat{R}$ couples position and heading through the lever arm to the tag. Tags farther
than \param{max_distance_m} (\SI{4}{m}) or seen beyond \param{max_view_angle_deg} (\ang{60}) are
ignored: at \SI{4}{m} a tag's position is already worth \SI{32}{cm}.

\section{The states}

\begin{description}
  \item[NONE] no $\T{map}{odom}$ yet: navigation does not start.
  \item[LOCALIZED] the robot knows where it is: on the start spot (\param{origin_sigma_xy}
    \SI{5}{cm}, \param{origin_sigma_yaw_deg} \ang{2}), or from a tag.
  \item[UNCERTAIN] $\sigma_{xy}$ above \param{max_sigma_xy} or $\sigma_\psi$ above
    \param{max_sigma_yaw_deg}.
\end{description}

The follower pauses only at NONE. UNCERTAIN is reported on \topic{localization/status} and by
the health monitor, and the patrol goes on: without a map the robot has nothing better than its
odometry between tags. The defaults (\SI{100}{m}, \ang{180}) never leave LOCALIZED; lower them to
make a long walk without tags visible in NETRA and the health report.

\section{Tags}

\begin{description}
  \item[Recording] \srv{localization/record_tags} averages each visible tag over
    \param{record_duration_sec} (\SI{2}{s}) and keeps it only with \param{record_min_samples}
    detections within \param{record_max_std_m} and \param{record_max_std_deg}. Record from where
    the robot will see the tag on its routes, \SIrange{1}{2.5}{m} away, as square-on as possible.
  \item[Automatic] \param{auto_localize}: while not localized, any visible recorded tag
    localizes the robot. \param{auto_correct}: while localized, visible tags are fused through the
    gate.
  \item[Placement] one tag every \SIrange{20}{40}{m} of route, at the camera's height, where the
    robot looks at it while walking (the start of a straight, a stop point), never on something
    that moves (a gate, a door).
\end{description}

\begin{warnbox}[A restarted odometry]
FAST-LIO restarting (Chapter~\ref{ch:t-odometry}) resets the odom frame: the old
$\T{map}{odom}$ means nothing. The localization drops to NONE and the patrol pauses. Show the
robot a recorded tag (\srv{localization/start}, the app's \ui{Localize}), or bring it back to its
start spot and restart the navigation host's service: the start spot becomes the origin again.
(\srv{localization/assume_origin} takes the current odom frame as the map: right only when
FAST-LIO restarted on the start spot.)
\end{warnbox}
